\section{Related Work}
\label{sec:related}

\subsection{Old Film and Video Restoration}
\label{subsec:ofr}

Old film restoration differs from general video restoration because it must jointly address structured artifacts, photometric degradation, and quality degeneration. Scratches, dust, and blotches require defect-aware repair, while flicker, grain, fading, blur, and noise require photometric and temporal correction. Traditional approaches decompose the problem into detection and inpainting stages, using hand-crafted features with interpolation-based filling~\cite{kokaram2004missing} or sparse collaborative filtering~\cite{dabov2007bm3d,he2013guided}. Deep learning shifted OFR toward end-to-end restoration: DeepRemaster~\cite{iizuka2019deepremaster} processes spatiotemporal volumes with 3D convolutions, RTN~\cite{wan2022rtn} introduces bidirectional recurrent propagation with a learnable soft defect mask, RRTN~\cite{lin2024rrtn} strengthens alignment with recursive transformer networks and second-order grid propagation, and MambaOFR~\cite{mao2025mambaofr} uses MambaIR-style~\cite{guo2024mambair} state-space modules while introducing the SRWOV benchmark.

Existing video restoration and enhancement systems provide useful context but are not tailored to OFR-specific defect structure. General video restoration methods such as EDVR~\cite{wang2019edvr}, BasicVSR~\cite{chan2021basicvsr}, BasicVSR++~\cite{chan2022basicvsrpp}, VRT~\cite{liang2024vrt}, and RVRT~\cite{liang2022rvrt} model temporal information effectively, but they do not explicitly handle old-film defects. Diffusion and foundation-prior systems, including DiffBIR~\cite{lin2024diffbir}, SUPIR~\cite{yu2024supir}, and SeeSR~\cite{wu2024seesr}, improve perceptual detail for blind image restoration or real-world super-resolution, but they are not designed around temporal old-film artifacts or SRWOV-style video protocols. This leaves a specific need for propagation features that can separate structural completion, regional repair, and fine artifact suppression.

\subsection{State Space Models for Vision}
\label{subsec:ssm}

State space models offer an efficient alternative to attention-heavy restoration backbones. The computational cost of self-attention scales quadratically with sequence length, which is costly for high-resolution video. Structured state space sequence models~\cite{gu2022s4} reformulate sequence modeling through continuous-time state transitions, reducing complexity from $O(L^2)$ to $O(L)$ while maintaining a global receptive field. Mamba~\cite{gu2024mamba} further introduces selective state transitions with hardware-aware recurrence. In vision, MambaIR~\cite{guo2024mambair} demonstrates 2D selective scanning for image restoration, VMamba~\cite{liu2024vmamba} builds hierarchical feature pyramids with four scanning routes, and VSRM~\cite{tran2025vsrm} alternates spatial and temporal Mamba blocks for video super-resolution at linear cost.

These properties are well suited to OFR, which combines long temporal dependency, large spatial damage, and high-resolution input. MambaOFR~\cite{mao2025mambaofr} shows that Mamba-based propagation can improve both efficiency and restoration quality relative to transformer counterparts. However, long-range aggregation alone does not specify which degradation scale each feature should represent. MgfrOFR therefore augments the SSM feature flow with explicit multi-granularity descriptors rather than replacing the backbone.

\subsection{Multi-Scale and Multi-Branch Representations}
\label{subsec:multigran}

Multi-scale and multi-branch representations are useful when one feature scale cannot capture both global context and fine local evidence. Feature Pyramid Networks~\cite{lin2017fpn} and HRNet~\cite{wang2021hrnet} show the value of explicit multi-scale decomposition for dense prediction, while MIRNet-v2~\cite{zamir2022mirnetv2} shows that complementary multi-scale context can improve low-level reconstruction. In person re-identification, PCB~\cite{sun2018pcb} decomposes features along spatial stripes, MGN~\cite{wang2018mgn} combines global and local branches at different partition granularities, and ReIDMamba~\cite{gu2025reidmamba} adapts this idea to state space models. These studies motivate branch-structured representations, but their objectives and feature semantics target recognition rather than video restoration.

\subsection{Diversity and Ranking Regularization}
\label{subsec:diversity_ranking}

Branch structure is useful only when the branches capture complementary cues rather than correlated copies of the same evidence. Contrastive learning~\cite{chen2020mocov2} provides a general diversity mechanism, but it does not directly regulate pairwise ordering behavior across branches. Two branches may differ in magnitude while ranking all candidate frames in the same order, producing apparent feature diversity but redundant temporal structure. Differentiable Kendall's Tau~\cite{zheng2023diffkendall} enables gradient-based penalization of correlated rankings, which is directly applicable to flow-guided temporal correspondences in OFR.

\subsection{Representation Symmetry in Restoration}
\label{subsec:symmetry_ai}

For this paper, representation symmetry denotes balanced and comparable organization of complementary restoration evidence rather than visible geometric symmetry. A learned representation becomes unbalanced when optimization favors one dominant cue and suppresses other cues needed for restoration. Feature-fusion and correlation-based representation studies show that multi-view or multi-branch features require explicit organization rather than simple concatenation~\cite{yu2021multiview_tensor_fusion,qian2020multiview_matrix_correlation}, while adjacent image-processing studies use attention, detection, or differential-equation priors to organize heterogeneous visual evidence for anomaly detection, object detection, segmentation, and denoising~\cite{cui2024video_anomaly_iqa_attention,yin2024improved_faster_rcnn,tian2020pde_segmentation_denoising}. Feature-wise modulation provides a practical way to inject conditioning signals into image-processing networks~\cite{perez2018film}. This restricted view is relevant to OFR because film damage is uneven across both spatial regions and frames.

MgfrOFR adapts these ideas to restoration rather than recognition. First, branch tokens are inserted into propagated restoration features rather than identity-oriented embeddings. Second, branch summaries are fed back into the restoration stream through channel-wise recalibration, so they affect reconstruction rather than only auxiliary supervision. Third, positives and negatives are defined by temporal correspondence within and across clips rather than person identities. The resulting design addresses two OFR-specific gaps: the lack of explicit multi-granularity decomposition within the temporal propagation pipeline, and the lack of diversity regularization that keeps parallel restoration branches complementary, as detailed in Section~\ref{sec:method}.
